\SourceLeaf{167}
Eher könnte es als zulässig erscheinen, für die Witwen zweiter Klasse neben der gleichmäßigen Beteiligung an der Zusatzdividende den festen Teil der Pension niedriger als 200 Rthl., zum Beispiel auf 150 Rthl., festzusetzen. Bedenkt man jedoch,

dass die Einbuße an Beiträgen sofort beginnt, sobald die zweite Klasse gebildet ist, und dann mit jeder Ernennung weiterer außerordentlicher Professoren fortwährend zunimmt;

dass sich jetzt unter den 51 Mitgliedern der Witwenkasse 20 außerordentliche Professoren befinden und folglich, wenn auch künftig ein ähnliches Verhältnis bestehen bleibt, von der im Hauptplan vorausgesetzten Beitragserhöhung ein Teil ausbleibt, der nach und nach bis auf \(\tfrac{2}{5}\) anwächst; während dagegen

der Ersatz durch die künftig anfallenden niedrigeren Pensionen aller Wahrscheinlichkeit nach erst nach langer Zeit beginnen und bei der sehr geringen Sterblichkeit in dem Alter, in dem die Mehrzahl der außerordentlichen Professoren gewöhnlich steht, durchschnittlich auch selten sein wird,

so bleibt es sehr fraglich, ob die Festsetzung des festen Pensionsteils auf 150 Rthl. für die zweite Klasse ausreichen würde, der Kasse auch nur den vollen Ersatz für den Verlust durch die niedrigeren Beiträge zu gewähren. Noch viel weniger dürfte man aber annehmen, eine solche Maßnahme gereiche der Kasse entschieden zum \textit{Vorteil}. [Dazu wäre man nur berechtigt, wenn man entweder die Beiträge in der zweiten Klasse ebenso hoch wie in der ersten beließe oder die Pension noch erheblich unter 150 Rthl. senkte. Beides würde jedoch so sehr als Härte erscheinen, dass sich die Kommission nicht dafür aussprechen kann.]

Eine andere Maßnahme, durch die der Kasse bisweilen ein gewisser Vorteil entstehen könnte, wäre die Aufhebung der Bestimmung im 10. Artikel des Regulativs, nach der die Witwenpension bei einer Wiederverheiratung der Witwe vollständig erlischt. Der Zweck dieser Vorschrift kann nur gewesen sein, der Kasse einen Gewinn zu verschaffen, weil man annahm, solche Fälle würden öfter vorkommen. Dieser Zweck wird aber so gut wie ganz verfehlt, da Wiederverheiratungen der Witwen unter diesen Umständen erfahrungsgemäß fast unerhört sind. (Schon in der Denkschrift ist bemerkt worden, dass in mehr als 100 Jahren nur ein solcher Fall vorgekommen ist.) Zweifellos zweckmäßiger ist die bei anderen Witwenkassen bestehende Einrichtung, dass die Witwenpension während einer zweiten Ehe nur ruht, aber wieder auflebt, wenn die Wiederverheiratete zum zweiten Mal Witwe wird. Man vergleiche die Statuten der Hof- und Civil-Diener Witwenkasse §. 25, der Prediger-Witwenkasse §. 24 und der Schullehrer-Witwenkasse §. 22. Der Fall, dass der zweite Ehemann wieder ein hiesiger Professor ist, müsste wohl ausgenommen werden, da es unzulässig scheint, einer Witwe (oder einer Familie) aus unserer Witwenkasse zwei Pensionen zu gewähren. Dagegen dürfte es ratsam sein, auf weitere Beschränkungen des Wiederauflebens zu verzichten und namentlich die volle Professorenwitwenpension auch dann wieder zu gewähren, wenn die Witwe durch ihre zweite Heirat, etwa mit einem Prediger oder anderen Staatsbediensteten, eine Pension aus einer \textit{anderen} öffentlichen Kasse erhielte. Man muss nämlich bedenken, dass von einer Änderung des bisherigen Zustands nur insoweit eine Wirkung erwartet werden kann, als die Änderung nicht durch Ausnahmebestimmungen wieder aufgehoben wird. Dass bei Kindern unter 20 Jahren aus der ersten Ehe die Waisenpension auch während der zweiten Ehe im bisherigen Umfang fortdauern müsste, versteht sich im Übrigen von selbst.

\Needspace{5\baselineskip}
\pdfbookmark[2]{Anmerkung zum Kommissionsbericht.}{witwen-iii-note}
\begin{center}\textit{Anmerkung zum Kommissionsbericht.}\end{center}
\noindent\hspace*{\parindent}Werden die Ausdrücke Bilanz und Defizit im Zusammenhang mit dem Finanzhaushalt eines Staates gebraucht, so versteht man unter Ersterer den Vergleich der Ausgaben und Einnahmen, wie sie für ein Jahr oder für eine kleine Anzahl von Jahren, die eine Finanzperiode bilden, nach voraussichtlicher Veran%
\SourceLeaf{168}schlagung erwartet werden. Unter Defizit versteht man den Unterschied, wenn sich für die Ausgaben eine größere Summe als für die Einnahmen ergibt. In einem solchen Zusammenhang sind Überschüsse und Defizit einander unmittelbar entgegengesetzt, und das eine schließt das andere von selbst aus.

Von einer solchen Rechnung unterscheidet sich die für die Witwenkasse in der zweiten Abteilung der Denkschrift durchgeführte Rechnung in zwei wesentlichen Punkten.

1) Letztere vergleicht nicht Ausgaben und Einnahmen für ein Jahr oder für einige Jahre, sondern umfasst beide für die ganze unbegrenzte Zukunft, soweit und so genau, wie es möglich ist, diese im Voraus zu veranschlagen.

2) Sie summiert nicht die Geldbeträge selbst in der Höhe, in der sie ausgegeben oder eingenommen werden (was wegen der Unbegrenztheit auch sinnlos wäre), sondern ihre Werte, die nach einem bestimmten Zinssatz auf den Anfangszeitpunkt abgezinst oder zurückgeführt sind. In diesem Sinn lässt sich auch eine unbegrenzt fortlaufende Reihe von Ausgaben oder Einnahmen zu einem bestimmten endlichen Ergebnis summieren. Diese Möglichkeit leuchtet leicht ein, wenn man bedenkt, dass eine Geldsumme durch Abzinsung über einen sehr langen Zeitraum ganz außerordentlich zusammenschrumpft: Eine nach 422 Jahren zu leistende Zahlung von 7000 Rthl. hat beispielsweise bei einem Zinssatz von \(3\tfrac{1}{2}\) Prozent jetzt nur den Wert von \textit{einem Pfennig}. Noch leichter gelangt man zum selben Ergebnis, wenn man bedenkt, dass der jetzige Geldwert einer jährlich ohne Ende in gleicher Höhe zu leistenden Zahlung nichts anderes ist als die Kapitalsumme, die bei dem gewählten Zinssatz jährlich einen ebenso großen Zins abwirft.

Nun ist zwar schon von selbst klar: Wenn nach den tatsächlichen Grundlagen einer solchen Rechnung die Ausgaben später größer sein werden als jetzt, die neben den Kapitalzinsen noch bestehenden Einnahmen in der Rechnung dagegen nicht als steigend angenommen werden dürfen, dann ergibt die Rechnung ein Defizit, falls jetzt die Einnahmen gegenüber den Ausgaben keinen Überschuss ergeben. Umgekehrt lässt sich aber nicht ohne Weiteres behaupten, dass in der Gesamtrechnung \textit{kein} Defizit bestehen werde, wenn jetzt die Einnahmen gegenüber den Ausgaben Überschüsse ergeben. Dazu sind nämlich nicht nur Überschüsse überhaupt, sondern Überschüsse von hinreichender Größe erforderlich. Hat eine richtige Gesamtbilanzrechnung als Endergebnis ein Defizit ergeben, so steht damit fest, dass die jetzigen jährlichen Überschüsse zu klein sind, um die künftig bevorstehenden Ausfälle zu decken. Die richtige Rechnung hat die gegenwärtigen zeitweiligen Überschüsse schon als Bestandteile der Bilanz mitberücksichtigt; sie dürfen der Anstalt nicht \textit{zweimal} als Einnahme angerechnet werden.

Ebenso falsch wie die Einbildung, dass jetzige zeitweilige Überschüsse ein jetzt bestehendes Defizit vermindern werden, ist die Vorstellung, dieses Defizit sei getilgt, sobald das Vermögen um den Betrag des Defizits gewachsen ist. Ein solcher Schluss ist nur in dem einzigen Fall zulässig, dass die Vermögensvermehrung \textit{sofort} erfolgt, und zwar durch \textit{fremde} Zuflüsse, die nicht schon in der Bilanz enthalten sind. Angenommen zum Beispiel, das Vermögen der Witwenkasse habe sich nach 20 Jahren (etwa teilweise durch neue Zuflüsse, die aber nicht fortdauernder Art wären) um 17520 Rthl. vermehrt: Dann wird das Defizit auch dann dennoch nicht beseitigt sein, sondern möglicherweise sogar größer sein als jetzt. Die Bilanz am 1. Oktober 1865 wird sich nämlich aus dem Vergleich des dann vorhandenen Kassenvermögens mit dem auf \textit{diesen} Zeitpunkt abgezinsten Betrag aller von da an bevorstehenden Ausgaben ergeben. Dieser Betrag wird viel größer sein als der entsprechende Betrag für den 1. Oktober 1845, und zwar aus dem einfachen Grund, dass 1865 die Zeit der viel höheren jährlichen Ausgaben um so viel näher gerückt sein wird. Ja, höchstwahrscheinlich werden diese dann bereits in beträchtlichem Umfang eingetreten sein.

